\textbf{Figure X | Trial-wise architecture of the dynamic Bayesian hypothesis model.}
\textbf{a}, A fixed labeled hypothesis space and participant-specific perceptual model provide the local/global transition kernels and rule-wise category predictions. The H dynamics use the previous trial's feedback surprise, posterior rule uncertainty and persistent control state to determine the revision mass $m_t$ and global-search share $g_t$. These controls redistribute probability over the full hypothesis space, producing the unique pre-feedback prior $\pi^-_t$ rather than hard-selecting a hypothesis. The same prior generates choice, RT and oral-report distributions. Feedback support then updates synchronized fade and static memory states, whose weighted combination gives $\pi^+_t$ and the inputs to the next trial. Solid arrows denote within-trial generation, dashed blue arrows denote next-trial recursion and dotted grey arrows denote fixed structural inputs. Stable maintenance, resistant maintenance, local revision and global search are descriptive regions of the continuous H state and do not control the generative process. \textbf{b}, Choice supplies the core behavioral constraint; RT and oral report are added sequentially to test whether they reduce predictive equivalence among memory, H and readout parameters. External validation freezes the choice-derived core posterior before testing the other channels, whereas joint integration quantifies contraction of the acceptable parameter set.
